%% Vol I, Chapter 3 — Physics-Driven Adaptive Runtime
%% SOURCE: docs/formal/scraps/architecture/03-architecture/physics-engine/feedback-loops.tex
%%         docs/formal/scraps/architecture/03-architecture/physics-engine/steady-state-machine.tex
%%         docs/formal/scraps/architecture/03-architecture/physics-engine/ssm-equations.tex
%%         docs/formal/scraps/architecture/03-architecture/physics-engine/metrics-and-knobs.tex
%%         docs/formal/scraps/architecture/03-architecture/adaptive-systems/adaptive-window-and-decay.tex
%%         docs/formal/scraps/architecture/03-architecture/pipelining/design.tex
%%         docs/formal/scraps/architecture/03-architecture/heartbeat-system/architecture.tex
%%         .claude/CLAUDE.md (Physics Engine section, Key Data Structures section)

\chapter{Physics-Driven Adaptive Runtime}
\label{vol1:chap:physics}

StarForth's adaptive runtime models its own execution behavior using execution
statistics as a proxy for thermodynamic quantities. Using execution frequency
as a proxy for thermal energy, the runtime tracks \emph{heat} per word, allows
heat to \emph{decay} over time, and reasons about \emph{stability} in
statistical terms. These are descriptive tools borrowed from thermodynamics as
a modeling convenience; they are not claims of physical thermodynamics. The
canonical term definitions are given in \texttt{ONTOLOGY.md}.

The runtime is governed by seven configurable feedback loops and an
overarching Steady-State Machine (SSM). Each loop pairs a positive
(accumulation) signal that drives the runtime toward optimization with a
negative (stabilization) signal that prevents runaway growth and holds the
system near equilibrium. All seven loops are independently togglable at build
time, enabling factorial Design of Experiments (DoE) campaigns that measure
each loop's contribution in isolation.

\section{Overview: Seven Feedback Loops}
\label{vol1:sec:physics:overview}

\begin{table}[ht]
\centering
\small
\begin{tabular}{llll}
\toprule
Loop & Name & Accumulation signal & Stabilization signal \\
\midrule
\#1 & Execution Heat       & word execution increments heat          & Loop \#3 decay; cold words demoted \\
\#2 & Rolling Window       & every execution recorded to buffer      & wrap overwrite; adaptive shrinking \\
\#3 & Linear Decay         & slope starts at $1/3$                   & decay reduces heat; slope adjusts $\pm 5\%$ \\
\#4 & Pipelining           & word$\rightarrow$word transitions recorded & counts age out; miss-rate feedback \\
\#5 & Window Inference     & window grows with diversity             & Levene's test shrinks to minimum \\
\#6 & Decay Inference      & regression extracts decay rate          & ANOVA early-exit; fit-quality bound \\
\#7 & Adaptive Heartrate   & tick counter increments                 & inference cadence falls when stable \\
\bottomrule
\end{tabular}
\caption{The seven feedback loops, with their accumulation and stabilization mechanisms.}
\label{vol1:tab:physics:loops}
\end{table}

Any loop can be disabled at build time by setting its flag to zero; setting
all seven to zero yields the baseline configuration used as the DoE reference
point. A global configuration table is given in
Section~\ref{vol1:sec:physics:knobs}.

\section{Loop \#1 --- Execution Heat}
\label{vol1:sec:physics:loop1}

%% SOURCE: feedback-loops.tex §Loop #1, .claude/CLAUDE.md §Physics Engine
%% Proof: proof/StarForth_Loop1_Heat.thy

Loop~\#1 counts word usage to identify hot execution paths. The implementation
resides in \texttt{src/dictionary\_heat\_optimization.c}. Every word execution
increments the \texttt{execution\_heat} field in the word's \texttt{DictEntry},
using execution frequency as a proxy for thermal energy. Words whose heat
crosses the promotion threshold migrate into the \texttt{hotwords\_cache} for
constant-time lookup; words falling below \texttt{HEAT\_CACHE\_DEMOTION\_THRESHOLD}
(default~10) are evicted from the cache. The \texttt{PHYSICS-RESET-STATS} word
resets all heat counters.

Two special flags modify decay behavior for individual words.
\texttt{WORD\_FROZEN} words never lose heat---their \texttt{execution\_heat}
does not participate in the Loop~\#3 decay sweep at all. \texttt{WORD\_PINNED}
words decay normally but are prevented from reaching zero: their heat is held at
a minimum of one after any decay step. These flags allow critical words to
remain permanently warm without requiring the heat tracking mechanism to be
suppressed globally.

Loop~\#1 interacts with Loop~\#3 (Linear Decay), which provides the
stabilization signal that prevents unlimited heat accumulation, and with the
hot-words cache (Section~\ref{vol1:sec:physics:hotwords}), which consumes
Loop~\#1's output. The build flag is \texttt{ENABLE\_LOOP\_1\_HEAT\_TRACKING}.
Formal properties of the heat lifecycle are established in
\texttt{proof/StarForth\_Loop1\_Heat.thy}, which proves heat non-negativity,
strict monotonicity of increment, an upper bound of \texttt{MAX\_HEAT},
immunity to decay for \texttt{WORD\_FROZEN} words, and a heat floor of~1 for
\texttt{WORD\_PINNED} words after any decay step.

\begin{description}
  \item[\texttt{HEAT\_CACHE\_DEMOTION\_THRESHOLD}] Heat value below which a
    word is evicted from the hot-words cache. Default: 10.
  \item[\texttt{ENABLE\_LOOP\_1\_HEAT\_TRACKING}] Build flag; set to~0 to
    disable the entire heat-tracking subsystem.
  \item[\texttt{WORD\_FROZEN}] Word flag: heat does not decay.
  \item[\texttt{WORD\_PINNED}] Word flag: heat cannot decay to zero.
\end{description}

\section{Loop \#2 --- Rolling Window of Truth}
\label{vol1:sec:physics:loop2}

%% SOURCE: feedback-loops.tex §Loop #2, .claude/CLAUDE.md §Key Data Structures
%% Proof: proof/StarForth_Loop2_Window.thy

Loop~\#2 captures the execution sequence in a circular buffer to provide
representative data for the inference engine. The implementation resides in
\texttt{src/rolling\_window\_of\_truth.c}. The key data structure is
\texttt{RollingWindowOfTruth}, which is embedded in the \texttt{VM} struct and
provides a circular buffer of recently executed word identifiers, double-buffered
snapshots for safe concurrent access by the heartbeat thread, and an adaptive
effective window size field controlled by Loops~\#2 and~\#5.

The window grows to \texttt{ROLLING\_WINDOW\_SIZE} (default~4096) before
wrapping and becomes ``warm'' after~1024 executions. Only a warm window holds
representative data; inference operations check the \texttt{is\_warm} flag
before running. Stabilization is twofold: oldest entries are overwritten on
wrap, and adaptive shrinking reduces \texttt{effective\_window\_size} whenever
pattern-diversity growth falls below \texttt{ADAPTIVE\_GROWTH\_THRESHOLD}
(1\%). The diversity check runs every \texttt{ADAPTIVE\_CHECK\_FREQUENCY}
(256) executions, retains \texttt{ADAPTIVE\_SHRINK\_RATE} (75\%) of the
current window per cycle, and never shrinks below
\texttt{ADAPTIVE\_MIN\_WINDOW\_SIZE} (256).

Loop~\#2 interacts with Loop~\#5, which applies Levene's test to determine
whether the window has reached its minimum sufficient size for the current
workload. The build flag is \texttt{ENABLE\_LOOP\_2\_ROLLING\_WINDOW}. Formal
properties are established in \texttt{proof/StarForth\_Loop2\_Window.thy},
which models both the effective window size and the actual (fill-level) window
size and proves that the effective window remains within the
$[\texttt{ADAPTIVE\_MIN}, \texttt{ROLLING\_WINDOW\_SIZE}]$ bounds.

\begin{description}
  \item[\texttt{ROLLING\_WINDOW\_SIZE}] Maximum window size and initial
    capacity. Default: 4096.
  \item[\texttt{ADAPTIVE\_CHECK\_FREQUENCY}] Executions between diversity
    checks. Default: 256.
  \item[\texttt{ADAPTIVE\_GROWTH\_THRESHOLD}] Pattern-diversity growth rate
    (\%) below which shrinking triggers. Default: 1\%.
  \item[\texttt{ADAPTIVE\_SHRINK\_RATE}] Fraction of window retained per
    shrink cycle. Default: 75\%.
  \item[\texttt{ADAPTIVE\_MIN\_WINDOW\_SIZE}] Floor for the effective window
    size. Default: 256.
\end{description}

\section{Loop \#3 --- Linear Decay}
\label{vol1:sec:physics:loop3}

%% SOURCE: feedback-loops.tex §Loop #3, adaptive-window-and-decay.tex
%% Proof: proof/StarForth_Loop3_Decay.thy

Loop~\#3 ages words to prevent unbounded heat accumulation; it is primarily a
negative-feedback loop. It does not have a separate implementation file but is
part of the heat-tracking subsystem in
\texttt{src/dictionary\_heat\_optimization.c}, driven by the heartbeat
dispatcher in \texttt{src/vm.c}. On each decay sweep, every word's heat is
reduced according to:

\begin{equation}
  H(t) = \max\!\bigl(0,\; H_0 - d \cdot \Delta t\bigr),
\end{equation}

where $d$ is the decay constant expressed in \Qtype{} fixed point as
\texttt{DECAY\_RATE\_PER\_US\_Q16}, giving a half-life of roughly six to seven
seconds for a word carrying 100 heat units. A minimum interval of
\texttt{DECAY\_MIN\_INTERVAL} ($1\,\mu$s) guards against over-decay when the
heartbeat fires in rapid succession.

The decay slope is stored as a \Qtype{} value (\texttt{decay\_slope\_q48},
initialized to~$1/3$) and self-adjusts by~$\pm 5$\% against the hot-to-stale
word ratio observed during each heartbeat sweep: an excess of hot words
increases the slope (faster decay), while an excess of stale words decreases
it (slower decay). This self-adjustment is the accumulation signal for a
primarily stabilizing loop. Loop~\#6 provides the longer-horizon inference
mechanism that recalculates the decay slope from exponential regression over
the heat trajectory.

The build flag is \texttt{ENABLE\_LOOP\_3\_LINEAR\_DECAY}. Formal properties
are established in \texttt{proof/StarForth\_Loop3\_Decay.thy}, which proves
that the decay operation is non-negative, monotonically decreasing, and that
\texttt{WORD\_FROZEN} words are immune to decay steps.

\begin{description}
  \item[\texttt{DECAY\_RATE\_PER\_US\_Q16}] Heat decay per microsecond in
    \Qtype{} fixed point. Default: $1/65536$ heat per $\mu$s.
  \item[\texttt{DECAY\_MIN\_INTERVAL}] Minimum elapsed time before a decay
    sweep fires. Default: $1\,\mu$s.
  \item[\texttt{decay\_slope\_q48}] Per-VM \Qtype{} value holding the current
    decay slope; initialized to~$1/3$; adjusted $\pm 5$\% by Loop~\#3's
    self-tuning and updated by Loop~\#6.
\end{description}

\section{Loop \#4 --- Word Transition Prediction (Pipelining)}
\label{vol1:sec:physics:loop4}

%% SOURCE: feedback-loops.tex §Loop #4, pipelining/design.tex
%%         metrics-and-knobs.tex §Pipelining Metrics
%% Proof: proof/StarForth_Loop4_Pipeline.thy

Loop~\#4 tracks word-to-word transitions to enable speculative prefetch,
reducing dictionary lookup latency by starting the lookup of word~$B$ during
the execution of word~$A$. The implementation resides in
\texttt{src/physics\_pipelining\_metrics.c}. Each \texttt{DictEntry} carries a
\texttt{WordTransitionMetrics} substructure that records transition counts
toward each successor, a total-transitions counter, and \Qtype{} latency
accounts for prefetch savings and misprediction cost. Global aggregate state
is held in \texttt{PipelineGlobalMetrics}, which includes
\texttt{prefetch\_attempts}, \texttt{prefetch\_hits},
\texttt{suggested\_next\_size} (the binary-chop window candidate), and
\texttt{last\_checked\_accuracy}. Transition probability is computed in \Qtype{}
fixed point to avoid floating-point arithmetic on the hot path:

\begin{equation}
  P(\text{from} \rightarrow \text{to})
  = \frac{\mathtt{transition\_heat}[\text{to}]}{\mathtt{total\_transitions}}.
\end{equation}

The design proceeds in phases. Instrumentation (recording transitions) is
implemented and active. Prediction analysis and the actual speculative
prefetch are subsequent phases. Transition counts age out over time, and a
miss-rate signal drives binary-chop search over
\texttt{TRANSITION\_WINDOW\_SIZE} through
\texttt{PipelineGlobalMetrics.suggested\_next\_size}. The build flag is
\texttt{ENABLE\_LOOP\_4\_PIPELINING\_METRICS}. Formal safety properties are
established in \texttt{proof/StarForth\_Loop4\_Pipeline.thy}.

\begin{description}
  \item[\texttt{TRANSITION\_WINDOW\_SIZE}] Context depth (number of prior
    words) used for transition prediction. Default: 2.
  \item[\texttt{SPECULATION\_THRESHOLD}] Minimum transition probability
    required before a speculative prefetch is issued. Default: 0.50.
  \item[\texttt{MIN\_SAMPLES}] Minimum recorded transitions before the
    predictor begins speculating. Default: 10.
  \item[\texttt{MISPREDICTION\_COST}] Penalty charged for a wrong prediction,
    in nanoseconds. Default: 25\,ns.
\end{description}

\section{Loop \#5 --- Window Width Inference}
\label{vol1:sec:physics:loop5}

%% SOURCE: feedback-loops.tex §Loop #5, adaptive-window-and-decay.tex §Toward a
%%         Bidirectional Loop
%% Proof: proof/StarForth_Loop5_WinInf.thy

Loop~\#5 finds the optimal rolling-window size through statistical testing, using
the inference engine in \texttt{src/inference\_engine.c}. The effective window
size expands toward \texttt{ROLLING\_WINDOW\_SIZE} when pattern diversity rises
above threshold. For stabilization, it applies Levene's test for equality of
variance: the heat trajectory is split into $K$ disjoint chunks, each chunk's
variance is computed independently, and the Levene test statistic $W$ is compared
against a critical value ($W \approx 6.5$ at $\alpha = 0.05$). When
$W \le$ critical, the variances are stable and the window has reached its minimum
sufficient size, preventing over-capture of redundant pattern data.

The binary-chop mechanism (\texttt{loop\_5\_binary\_chop\_suggest\_window()})
is invoked by the heartbeat window-tuner plugin
(\texttt{vm\_tick\_window\_tuner()}, which runs at
\texttt{WINDOW\_TUNING\_FREQUENCY} ticks). It computes the current prefetch
accuracy and proposes a new effective window size, which the window tuner
applies if it differs from the current size.

The Levene test's statistical correctness---that it correctly detects
non-uniform execution diversity---is taken as an explicit named axiom
(\texttt{inference\_window\_clamped\_valid}) in
\texttt{proof/StarForth\_Loop5\_WinInf.thy}, carrying a documented human-audit
obligation against \texttt{src/inference\_engine.c}. All other properties
(boundedness, monotonicity, ANOVA early-exit) are fully mechanised.

\begin{description}
  \item[\texttt{WINDOW\_TUNING\_FREQUENCY}] Heartbeat ticks between
    window-tuner invocations. Default: 1000.
  \item[\texttt{ADAPTIVE\_MIN\_WINDOW\_SIZE}] Hard floor on the effective
    window size. Default: 256.
\end{description}

\section{Loop \#6 --- Decay Slope Inference}
\label{vol1:sec:physics:loop6}

%% SOURCE: feedback-loops.tex §Loop #6
%% Proof: proof/StarForth_Loop6_DecayInf.thy

Loop~\#6 extracts the optimal decay rate by exponential regression over the
heat trajectory, using the inference engine in \texttt{src/inference\_engine.c}.
Fitting the log-linear model

\begin{equation}
  \ln(\mathrm{heat}[t]) = \ln(h_0) - \mathrm{slope} \cdot t
\end{equation}

yields a closed-form \texttt{adaptive\_decay\_slope}. An ANOVA early-exit
guards cost: if \texttt{has\_variance\_stabilized()} reports that variance
changed by less than~5\% since the last inference, the computation is skipped
and the cached result reused, saving an estimated 5{,}000--10{,}000 CPU cycles
per skipped run. The slope is bounded by fit quality
(\texttt{slope\_fit\_quality\_q48}), and
\texttt{inference\_outputs\_validate()} rejects slopes that are zero or exceed
100.

The build flag is \texttt{ENABLE\_LOOP\_6\_DECAY\_INFERENCE}. Unlike
Loop~\#5, Loop~\#6 requires no statistical axiom: all invariants are proved
purely from the clamping of the raw ordinary least-squares slope to
$[\texttt{DECAY\_SLOPE\_MIN}, \texttt{DECAY\_SLOPE\_MAX}]$ before writing to
the VM state. This is established in \texttt{proof/StarForth\_Loop6\_DecayInf.thy}.

\begin{description}
  \item[\texttt{slope\_fit\_quality\_q48}] \Qtype{} quality metric for the
    current regression fit; used to gate slope updates.
  \item[\texttt{HEARTBEAT\_INFERENCE\_FREQUENCY}] Ticks between full inference
    runs. Default: 5000.
  \item[\texttt{SLOPE\_VALIDATION\_FREQUENCY}] Ticks between decay-slope
    validation passes. Default: 5000.
\end{description}

\section{Loop \#7 --- Adaptive Heartrate}
\label{vol1:sec:physics:loop7}

%% SOURCE: feedback-loops.tex §Loop #7, heartbeat-system/architecture.tex
%%         .claude/CLAUDE.md §Critical Implementation Details
%% Proof: proof/StarForth_Loop7_Heartrate.thy

Loop~\#7 adjusts tick frequency to balance responsiveness against inference
cost. All time-driven VM operations are coordinated through a single
dispatcher, \texttt{vm\_tick()}, in \texttt{src/vm.c}. The \texttt{HeartbeatState}
struct, embedded in the \texttt{VM} struct and declared in
\texttt{include/vm.h}, holds the tick counter and per-plugin cursors:

\begin{lstlisting}[language=C]
typedef struct {
    uint64_t tick_count;            /* total ticks since VM init */
    uint64_t last_window_tune_tick; /* tick of last window tune */
    uint64_t last_slope_tune_tick;  /* tick of last decay validation */
    int      heartbeat_enabled;     /* 1 = active, 0 = disabled */
} HeartbeatState;
\end{lstlisting}

The main interpreter increments an execution counter and fires
\texttt{vm\_tick()} every \texttt{HEARTBEAT\_FREQUENCY} (default~256) word
executions. The dispatcher increments \texttt{tick\_count} and invokes each
registered plugin at its configured interval: the window-tuner plugin
(\texttt{vm\_tick\_window\_tuner()}, every \texttt{WINDOW\_TUNING\_FREQUENCY}
ticks) and the slope-validator plugin (\texttt{vm\_tick\_slope\_validator()},
every \texttt{SLOPE\_VALIDATION\_FREQUENCY} ticks).

When the system is stable, full inference runs less often, reducing CPU
overhead. The background \texttt{HeartbeatWorker} thread can be enabled with
\texttt{HEARTBEAT\_THREAD\_ENABLED=1}, in which case it wakes at
\texttt{HEARTBEAT\_THREAD\_INTERVAL\_MS} and calls \texttt{vm\_tick()}
independently of the execution hot-path. In threaded mode, state shared between
the execution thread and the heartbeat thread is protected by
\texttt{vm->tuning\_lock}.

\textbf{Implementation status:} \texttt{HeartbeatTickSnapshot} and
\texttt{tick\_buffer} are declared in \texttt{include/vm.h}, but
\texttt{heartbeat\_export\_csv()} is not yet implemented. The instrumentation
plan is recorded in
\texttt{docs/03-architecture/heartbeat-system/instrumentation-plan.md}.

The build flag is \texttt{ENABLE\_LOOP\_7\_ADAPTIVE\_HEARTRATE}. Formal
invariants are established in \texttt{proof/StarForth\_Loop7\_Heartrate.thy},
which proves that \texttt{tick\_target\_ns} is always greater than zero,
remains within $[\texttt{TICK\_MIN\_NS}, \texttt{TICK\_MAX\_NS}]$, and that
\texttt{hb\_tick\_count} increases monotonically.

\begin{description}
  \item[\texttt{HEARTBEAT\_FREQUENCY}] Word executions between ticks.
    Default: 256.
  \item[\texttt{WINDOW\_TUNING\_FREQUENCY}] Ticks between window-tuner
    calls. Default: 1000.
  \item[\texttt{SLOPE\_VALIDATION\_FREQUENCY}] Ticks between decay-slope
    validation. Default: 5000.
  \item[\texttt{HEARTBEAT\_THREAD\_ENABLED}] Build flag; 1 enables the
    background POSIX thread. Default: 0 (synchronous mode).
  \item[\texttt{HEARTBEAT\_THREAD\_INTERVAL\_MS}] Wake interval for the
    background thread when enabled. Default: 100\,ms.
\end{description}

\section{L8 Jacquard Mode Selector}
\label{vol1:sec:physics:l8}

%% SOURCE: .claude/CLAUDE.md §L8 Jacquard Mode Selector, ssm_jacquard.c
%%         steady-state-machine.tex

The L8 Jacquard Mode Selector is a meta-coordinator that sits above the seven
feedback loops. It is implemented in \texttt{src/ssm\_jacquard.c} and its
current mode is held in \texttt{vm->ssm\_l8\_state}. Rather than a feedback
loop in the same sense as Loops~\#1--\#7, L8 is a Steady-State Machine (SSM)
that switches the runtime between predefined operating configurations based on
observed attractor-bucket statistics aggregated by the heartbeat mechanism.

Using execution statistics as a proxy for thermodynamic quantities, the SSM
tracks execution heat, entropy (workload diversity), pipeline readiness, heat
decay, and the steady-state slope. When a configuration's attractor-bucket
counts reach characteristic thresholds, L8 transitions the runtime to a
different operating mode, adjusting the balance of loops and their parameters
to suit the new workload regime. The SSM passes through four phases: transient
(warming up, noisy gradients), quasi-steady (emerging patterns), true steady
state (smooth heat surface, stable lookups), and chaotic/disruption
(triggered by major workload shift).

Six operating modes are defined:

\begin{description}
  \item[stable] The system has converged to a predictable workload; decay
    rates and window sizes are held conservatively to avoid disrupting an
    established optimum.
  \item[volatile] High workload diversity; the window is widened and
    inference runs more frequently to track rapidly changing heat surfaces.
  \item[diverse] Moderate, sustained diversity; window and decay parameters
    are tuned for broad pattern capture.
  \item[temporal] Time-structured workloads where periodically recurring
    patterns dominate; the window is sized to capture at least one full
    period.
  \item[transition] The system is moving between attractors; parameters are
    relaxed to allow rapid re-convergence.
  \item[omni] All-loops active at maximum sensitivity; intended for initial
    warm-up or after a major disruption.
\end{description}

The L8 mode selector is not a loop with an accumulation/stabilization pair;
it is a higher-order coordinator that reconfigures the parameters of the
seven loops based on observed steady-state behavior. The six L8 variants are
also available as separate capsule personalities:
\texttt{capsules/init-l8-\{stable,volatile,diverse,temporal,transition,omni\}.4th}.

\section{Steady-State Machine Theory}
\label{vol1:sec:physics:ssm}

%% SOURCE: steady-state-machine.tex

The theoretical basis for the adaptive runtime is the Steady-State Machine
(SSM), a computational model whose operational state is defined not by discrete
symbolic transitions but by convergence to an optimal runtime equilibrium.
Using execution statistics as a proxy for thermodynamic quantities, the SSM
tracks execution heat $H$, entropy and diversity measures $E$, pipeline
activation thresholds $P$, a heat-decay profile $\Theta$, and the sliding
execution window $W$. The governing dynamic is:

\begin{equation}
M_{t+1} = f\!\bigl(M_t,\; \Delta H,\; \Delta W,\; \Delta E\bigr),
\end{equation}

where $f$ drives the machine toward a stable basin. The machine is defined by
its convergence behavior, not by its instruction sequence.

Five axioms characterize the SSM:

\begin{description}
  \item[Convergence] Every sustained workload induces the SSM to converge
    toward a stable attractor unless external entropy forces divergence.
  \item[Locality of Heat] Execution heat reflects both locality and temporal
    relevance; locality produces stability.
  \item[Adaptive Reordering] Dictionary reordering is thermodynamic
    self-optimization, not symbolic mutation.
  \item[Stability Maximizes Throughput] The steady state is always faster
    than the cold state.
  \item[Perturbation Response] When disrupted, the SSM seeks a new steady
    state; it does not thrash indefinitely.
\end{description}

\section{Configuration Knobs}
\label{vol1:sec:physics:knobs}

%% SOURCE: feedback-loops.tex §Configuration

Any loop can be disabled at build time. Setting all seven flags to zero yields
the baseline configuration used as the DoE reference point.

\begin{lstlisting}[language=bash]
# Disable a single loop
make ENABLE_LOOP_3_LINEAR_DECAY=0

# Baseline build (all loops off)
make ENABLE_LOOP_1_HEAT_TRACKING=0 \
     ENABLE_LOOP_2_ROLLING_WINDOW=0 \
     ENABLE_LOOP_3_LINEAR_DECAY=0 \
     ENABLE_LOOP_4_PIPELINING_METRICS=0 \
     ENABLE_LOOP_5_WINDOW_INFERENCE=0 \
     ENABLE_LOOP_6_DECAY_INFERENCE=0 \
     ENABLE_LOOP_7_ADAPTIVE_HEARTRATE=0
\end{lstlisting}

\begin{table}[ht]
\centering
\small
\begin{tabular}{lll}
\toprule
Knob & Default & Description \\
\midrule
\texttt{ROLLING\_WINDOW\_SIZE}            & 4096 & Initial/maximum window size \\
\texttt{ADAPTIVE\_SHRINK\_RATE}           & 75   & \% retained when shrinking \\
\texttt{ADAPTIVE\_MIN\_WINDOW\_SIZE}      & 256  & Floor for window shrinking \\
\texttt{ADAPTIVE\_CHECK\_FREQUENCY}       & 256  & Executions between diversity checks \\
\texttt{ADAPTIVE\_GROWTH\_THRESHOLD}      & 1    & Growth rate (\%) signalling saturation \\
\texttt{DECAY\_RATE\_PER\_US\_Q16}        & 1    & Heat decay per microsecond (\Qtype{}) \\
\texttt{HEARTBEAT\_INFERENCE\_FREQUENCY}  & 5000 & Ticks between full inference runs \\
\texttt{TRANSITION\_WINDOW\_SIZE}         & 2    & Context depth for pipelining prediction \\
\bottomrule
\end{tabular}
\caption{Principal tuning knobs for the feedback loops.}
\label{vol1:tab:physics:knobs}
\end{table}

The implementation is distributed across
\texttt{src/rolling\_window\_of\_truth.c} (Loop~\#2),
\texttt{src/dictionary\_heat\_optimization.c} (Loops~\#1 and~\#3),
\texttt{src/inference\_engine.c} (unified Loops~\#5 and~\#6),
\texttt{src/physics\_pipelining\_metrics.c} (Loop~\#4), and the
\texttt{vm\_tick()} heartbeat dispatcher in \texttt{src/vm.c} (Loop~\#7).
Knob definitions live in \texttt{include/rolling\_window\_knobs.h}.

\section{Determinism and Formal Properties}
\label{vol1:sec:physics:determinism}

%% TODO(bob): 0\% algorithmic variance claim from the 90-run DoE.
%% Cite james:2025:ssrn and the relevant Isabelle proofs.
%% Source: docs/working/papers/NULL_HYPOTHESIS.md for falsification framing.

\section{Hot-Words Cache Optimization}
\label{vol1:sec:physics:hotwords}

%% SCRAP: experiments/02-experiments/physics-optimization/hotwords-cache
%% SOURCE: docs/working/experiments/02-experiments/physics-optimization/hotwords-cache.md
%% STATUS: HISTORICAL
%% FITS: experiments/ch-physics-opt, vol1-vm-physics/ch-loops
%% EDITORIAL: lifted — prose rewritten to press voice

\section{Hot-Words Cache Performance Analysis}
\label{sec:hotwords-cache-results}

\subsection{Summary}

The StarForth hot-words cache achieves a statistically confirmed
\textbf{1.78$\times$} speedup for dictionary lookups, with a 95\%
Bayesian credible interval of $[1.75\times, 1.81\times]$. All
measurements use Q48.16 fixed-point arithmetic; no floating-point
computation is required.

\subsection{Experimental Configuration}

\begin{itemize}
  \item \textbf{Build:} \texttt{make ENABLE\_HOTWORDS\_CACHE=1 fastest}
        (x86\_64, full assembly optimisations, LTO, direct threading)
  \item \textbf{Benchmark:} 100{,}000 dictionary lookups via FORTH test
        harness; word mix: common FORTH primitives (\texttt{IF}, \texttt{DUP},
        \texttt{DROP}, \texttt{+}, \texttt{-}, \texttt{@}, \texttt{!}, etc.)
  \item \textbf{Cache:} 32~entries; promotion threshold: 50~executions;
        eviction: LRU
  \item \textbf{Precision:} Q48.16 64-bit fixed-point (nanoseconds)
\end{itemize}

\subsection{Lookup Distribution}

\begin{center}
\begin{tabular}{lrr}
\toprule
Path & Count & Fraction \\
\midrule
Cache hits   & 1{,}415 & 35.64\% \\
Bucket hits  & 1{,}861 & 46.88\% \\
Misses       &   694   & 17.48\% \\
\midrule
Total        & 3{,}970 & 100\% \\
\bottomrule
\end{tabular}
\end{center}

\subsection{Latency Results}

\begin{center}
\begin{tabular}{lrrr}
\toprule
Path & Min (ns) & Mean (ns) & Max (ns) \\
\midrule
Cache path (1{,}415 samples)  & 22 & 31.543 & 101 \\
Bucket path (1{,}861 samples) & 23 & 56.237 & 370 \\
\bottomrule
\end{tabular}
\end{center}

\subsection{Speedup Analysis}

\begin{equation}
  \text{speedup} = \frac{\bar{t}_{\text{bucket}}}{\bar{t}_{\text{cache}}}
                 = \frac{56.237\,\text{ns}}{31.543\,\text{ns}}
                 = 1.78\times
\end{equation}

Time saved per cache hit: $56.237 - 31.543 = 24.694$~ns (43.9\% reduction).

\paragraph{Bayesian credible intervals.}

\begin{center}
\begin{tabular}{ll}
\toprule
Interval & Speedup \\
\midrule
95\% credible & $[1.75\times,\ 1.81\times]$ \\
99\% credible & $[1.73\times,\ 1.83\times]$ \\
\bottomrule
\end{tabular}
\end{center}

$P(\text{speedup} > 1.1\times) = 99.9\%$. All credible-interval arithmetic
is performed in pure Q48.16 integer arithmetic; no floating-point or libm
functions are used.

\subsection{Cache Management Behaviour}

The physics engine promoted 10~words automatically from the dictionary into
the cache based on execution entropy exceeding the threshold of~50. No
manual tuning was applied. Zero evictions occurred during the benchmark;
the 32-entry cache was sufficient for the active working set.

The two highest-entropy words at experiment completion were \texttt{EXIT}
(entropy~114) and \texttt{LIT} (entropy~101), consistent with their
frequency in the compiled FORTH test harness.

\subsection{Production Properties}

\begin{itemize}
  \item \textbf{No floating-point:} All arithmetic is Q48.16 integer;
        the cache subsystem is compatible with L4Re and bare-metal targets.
  \item \textbf{No external libraries:} No libm dependency.
  \item \textbf{Deterministic:} Near-zero variance across 3{,}970 samples
        confirms formally proven deterministic behaviour.
  \item \textbf{Linear complexity:} Cache lookup is O(1); bucket search
        is linear in bucket depth.
\end{itemize}

\subsection{Test Suite Validation}

\begin{lstlisting}[language=bash]
make fastest ENABLE_HOTWORDS_CACHE=1
make test
# Total: 782  Passed: 731  Failed: 0  Skipped: 49  Errors: 0
\end{lstlisting}

All 731 implemented tests pass with the cache-enabled build.



\section{Optimization Proposals}
\label{vol1:sec:physics:proposals}

%% SCRAP: experiments/02-experiments/physics-optimization/proposals
%% SOURCE: docs/working/experiments/02-experiments/physics-optimization/proposals.md
%% STATUS: CURRENT
%% FITS: experiments/ch-physics-opt, vol1-vm-physics/ch-loops
%% EDITORIAL: lifted — prose rewritten to press voice

\section{Physics-Driven Optimisation Proposals}
\label{sec:physics-opt-proposals}

\subsection{Strategic Direction}

The hot-words cache experiment (1.78$\times$, 95\% CI: [1.75$\times$,
1.81$\times$]) validates the physics-driven optimisation methodology:
collect execution-frequency metrics at runtime, make automatic promotion
decisions, and measure impact with statistical rigour. Nine additional
opportunities build on the same foundation.

JIT compilation is explicitly excluded from consideration. JIT requires
runtime code generation, violates L4Re microkernel policy, introduces
floating-point overhead, and defeats formal verification. Physics-driven
optimisation achieves measurable gains within StarForth's verification and
portability constraints.

\subsection{Opportunity Catalogue}

\paragraph{Opportunity~1 — Return Stack Prediction and Colon Word Inlining.}
Frequently called colon definitions with small bodies (fewer than 256~bytes)
are candidates for compile-time inlining when \texttt{execution\_heat > 100}.
Inlining eliminates dictionary lookup and return-stack overhead.
Expected gain: 1.3--2.0$\times$ for call-heavy workloads. Verification:
static (inlining correctness is provable via Isabelle/HOL).

\paragraph{Opportunity~2 — Block I/O Prefetching.}
A Markov-style transition matrix tracks which block LBN follows which.
When a block is accessed and the next-block heat exceeds a threshold, that
block is pre-fetched into the buffer cache. Expected gain: 1.5--3.0$\times$
for block-sequential access patterns. Implementation: medium complexity.

\paragraph{Opportunity~3 — Stack Operation Fusion.}
Co-execution heat tracks consecutive word pairs (e.g., \texttt{DUP DROP},
\texttt{SWAP ROT}, \texttt{OVER SWAP}). Pairs exceeding a heat threshold
at compile time are fused into dedicated primitives, eliminating two lookups
and one execution per pair. Expected gain: 1.2--1.5$\times$ for stack-heavy
programs. Verification: static (composition of verified primitives).

\paragraph{Opportunity~4 — Memory Allocation Pattern Prediction.}
Allocation frequency by size is tracked in a heat vector. Sizes exceeding
a threshold trigger pre-warmed pool maintenance. Repeating allocation
sequences are detected and cached. Expected gain: 1.3--1.8$\times$ for
allocation-heavy programs.

\paragraph{Opportunity~5 — Control Flow Branch Prediction.}
\texttt{IF}/\texttt{THEN}/\texttt{ELSE} outcomes are tracked per source
location. Branches with greater than 90\% one-sided bias are annotated
with a prediction hint that the inner interpreter can use to reorder code
or emit x86 branch-hint prefixes. Expected gain: 1.1--1.3$\times$
(architecture-dependent).

\paragraph{Opportunity~6 — Vocabulary Search Path Reordering.}
Per-vocabulary hit rates are tracked. The search order is sorted
dynamically by hit rate, placing the most productive vocabulary first.
Expected gain: 1.2--1.6$\times$ in multi-vocabulary programs.
Implementation: low complexity.

\paragraph{Opportunity~7 — String Operation Batching.}
Consecutive string-output operations (\texttt{."}, \texttt{EMIT},
\texttt{TYPE}) detected at compile time are fused into a single batch
output, reducing interpreter loop iterations. Expected gain:
1.1--1.4$\times$ for I/O-bound programs. Implementation: low complexity.

\paragraph{Opportunity~8 — Arithmetic Operation Reordering.}
For hot commutative operations at a given source location, operand sizes
are compared and the smaller operand is loaded first to improve prefetch
behaviour. Expected gain: 1.05--1.15$\times$ (architecture-dependent).

\paragraph{Opportunity~9 — Word Placement Optimisation.}
Co-execution heat between word pairs guides memory compaction: frequently
co-executed words are relocated adjacent to each other in the dictionary
to improve instruction-cache locality. Expected gain: 1.05--1.2$\times$.
Implementation: high complexity (requires GC integration).

\subsection{Implementation Roadmap}

\begin{center}
\begin{tabular}{lll}
\toprule
Phase & Opportunities & Rationale \\
\midrule
1 (complete) & Hot-words cache & Proven; production-ready \\
2 (recommended next) & \#3, \#6, \#1 & Highest ROI, lowest complexity \\
3 (future) & \#2, \#7, \#4 & Medium complexity \\
4 (advanced) & \#5, \#8, \#9 & CPU-specific or GC-dependent \\
\bottomrule
\end{tabular}
\end{center}

\subsection{Expected Cumulative Gains}

Conservative sequential composition of Phases~1--2:

\[
  1.78\times \;\times\; 1.2\;\times\; 1.2\;\times\; 1.3
  \;\approx\; 3.3\times \text{ cumulative speedup}
\]

All nine opportunities together: 5--8$\times$ total improvement is
plausible, subject to workload dependency and diminishing returns.

\subsection{Unified Metrics Infrastructure}

All nine opportunities require co-execution heat tracking and sequence
pattern buffers not yet present in the VM. A one-time extension to the
\texttt{PhysicsMetrics} per-word structure adds: a co-execution heat
vector, timing accumulator, per-word cache-line counters, and a small
recent-execution circular buffer. This single investment unlocks the
infrastructure required for Opportunities~1--9.


